\documentclass[a4paper,12pt]{article}

\usepackage{alltt, fancyvrb, url}
\usepackage{graphicx}
\usepackage[utf8]{inputenc}
\usepackage{float}
\usepackage{hyperref}
\usepackage[italian]{babel}
\usepackage[table]{xcolor}
\usepackage{array}
\usepackage{multirow}
\usepackage{titlesec}
\usepackage{listings}
\usepackage[italian]{cleveref}
\usepackage{acronym}
\usepackage{subcaption}

\renewcommand{\thesection}{\arabic{section}}

\sloppy
\setlength{\parindent}{0pt}

\title{\textbf{Chiusura del progetto}}
\author{}
\date{}

\begin{document}
\maketitle
\tableofcontents

\newpage

\section{Procedura di accettazione}
Il cliente effettua il collaudo della soluzione consegnata, verificando il rispetto dei criteri di accettazione
stabiliti ed approvati in fase di planning, oltre alle CoS stabilite nel POS. In caso positivo, il cliente certifica
l'accettazione della soluzione consegnata.

\section{Modalità di installazione}
L'installazione della soluzione avviene seguendo un \textit{Phased approach}, quindi verrà installata l'ultima parte
della soluzione, in quanto le prime parti sono già state consegnate ed installate in precedenza.

\section{Audit post-implementazione}
Si svolge l'audit post-implementazione, che consiste in un'analisi della soluzione consegnata e del processo di
sviluppo, in modo da effettuare una retrospettiva per migliorare in progetti futuri. \\
Un esempio di audit post-implementazione è il seguente:
\begin{itemize}
    \item gli obiettivi sono stati raggiunti e il deliverable era conforme con le aspettative del cliente;
    \item il progetto è stato completato rispettando i limiti di budget e tempo;
    \item il committente si è mostrato soddisfatto del risultato;
    \item Kanban era stato poco utilizzato in progetti passati, dunque è stata una possibilità per migliorare la
    conoscenza della metodologia. In particolare, sono stati notati i vantaggi legati alla flessibilità nelle durate
    dei cicli di sviluppo caratterizzanti Scrum (metodologia più utilizzata in passato).
\end{itemize} 

\end{document}